\section{Whose Image Is This?}
\label{sec:opening}

The Gospel of this Sunday is a trap and its undoing. Jesus is teaching in the
temple at Jerusalem, and the men he has been answering go away to plan.
``Then the Pharisees going, consulted among themselves how to insnare him in
his speech. And they sent to him their disciples with the Herodians, saying:
Master, we know that thou art a true speaker and teachest the way of God in
truth. Neither carest thou for any man: for thou dost not regard the person of
men. Tell us therefore what dost thou think? Is it lawful to give tribute to
Caesar, or not?'' Jesus, ``knowing their wickedness'', calls them hypocrites
and asks for the coin. ``Shew me the coin of the tribute. And they offered him
a penny. And Jesus saith to them: Whose image and inscription is this? They
say to him: Caesar's. Then he saith to them: Render therefore to Caesar the
things that are Caesar's; and to God, the things that are God's'' (Mt
22:15--21).

{\small\noindent Scripture is quoted throughout in the Douay--Rheims version as
revised by Bishop Challoner, a public-domain study translation made from the
Latin Vulgate. It is not the text proclaimed at Mass: the English heard in the
celebration is that of the \work{Lectionary for Mass for Use in the Dioceses of
the United States of America}, which alone is the approved wording. Psalms are
numbered as the Lectionary numbers them, with the Vulgate number, which the
Douay--Rheims and the Latin books use, in parentheses; where a psalm verse is
quoted from the Douay--Rheims, its Vulgate number is given.\par}
\medskip

The answer is famous and brief, and it leaves a great deal open. The coin
belongs to Caesar because his image is on it; that much the questioners say
themselves. But Jesus does not say what ``the things that are God's'' are, nor
how far Caesar's claim runs, nor why men who had praised him truthfully should
be called hypocrites. The rest of the Mass surrounds the scene with texts that
press on those questions from different sides.

The first reading is an oracle addressed to a foreign king. ``Thus saith the
Lord to my anointed Cyrus, whose right hand I have taken hold of, to subdue
nations before his face, and to turn the backs of kings.'' The Lord has called
the Persian by name for the sake of Israel, ``and thou hast not known me'';
he has girded him, so that from the rising of the sun to the west all may
know ``that there is none besides me'' (Is 45:1, 4--6). The psalm answers by
summoning the nations to bring the Lord the glory due to him, declaring that
the gods of the gentiles are nothing beside the Lord who made the heavens, and
announcing among them that the Lord reigns and will judge the peoples with
justice (Ps 96 (95)). The second reading is the opening of Paul's first letter
to the church of the Thessalonians, a church ``in God the Father and in the
Lord Jesus Christ'', whose work of faith, labour of charity and enduring hope
Paul remembers before God, and to whom the gospel came ``not \ldots\ in word
only, but in power also'' (1~Thes 1:1--5b). The Alleluia verse, taken from
Paul's letter to the Philippians, speaks of those who shine as lights in the
world, holding forth the word of life (Phil 2:15d, 16a).

The Missal's own texts, which serve this Sunday in all three years of the
Lectionary's cycle, frame the readings with three antiphons and three prayers.
The Entrance Antiphon is the cry of one who calls on God because God has heard,
and asks to be kept as the apple of the eye and sheltered under the shadow of
his wings. The Collect asks that the assembly always hold its will devoted to
God and serve his majesty with an undivided heart. The Prayer over the
Offerings asks for a free mind in the service of God's own gifts, and the
Prayer after Communion that frequent sharing in heavenly things may help the
assembly forward, with temporal benefits as its aid and eternal things as its
instruction. At Communion the Missal offers either the psalm verse on the eyes
of the Lord that rest on those who fear him and hope in his mercy, or the
saying that the Son of Man came to give his life as a ransom for many.

The \work{Ordo lectionum Missae}, on which the Lectionary rests, marks what it
singles out in each reading by the title it sets over it. Over the first
reading stands \latin{Apprehendi dexteram Cyri ut subiciam ante faciem eius
gentes}, from Isaiah 45:1; over the Gospel, \latin{Reddite quae sunt Caesaris
Caesari, et quae sunt Dei Deo}, from Matthew 22:21; over the second reading,
\latin{Memores fidei vestrae et caritatis et spei}, from 1~Thessalonians 1:3.
The introduction to the \work{Ordo} explains that the relation between the
readings of one Mass is shown by the careful choice of these titles, and that
in Ordinary Time the Old Testament reading is chosen for its relation to the
Gospel while the Apostle and the Gospel are each read in their own course
(\work{Praenotanda} 105--107). The two titles that belong together each name a
gentile ruler: Cyrus, whose right hand the Lord grasps, and Caesar, who is to
be given what is his. The second reading opens the Lectionary's course through
the First Letter to the Thessalonians in Year~A, and the Gospel continues its
course through Matthew.

The Sunday's question, then, is what is owed, and to whom, and by what kind of
giver. The Fathers and Doctors, each expounding his own passage, answer it
along three lines, and each line, carried through every text of the Mass,
yields an interpretation of the whole.

The first interpretation asks what ``the things that are God's'' are. It
follows an inference from the coin's image to the human person, made in the
image of God. Augustine introduces that inference as what ``some understand
skilfully'' and then makes it his own; Ambrose, on Luke's form of the saying,
says that Christ bears no image of Caesar because he is the image of God;
Hilary of Poitiers reaches the debt of the whole self by another premise, that
body, soul and will are owed to God because they come from him; Thomas Aquinas
gives a literal sense and keeps the image reading as the mystical one. Jerome
reads the second half of the saying literally, of tithes and offerings. The
second asks how God's people can owe
anything to a ruler who does not know God, and how far the debt runs. It reads
the Gospel beside the oracle to Cyrus, whom Jerome and Aquinas explain as a
gentile king called the Lord's anointed for his royal dignity; it takes
from John Chrysostom that tribute is rendered as a debt and ends where
godliness would be harmed, and from Augustine the soldiers who served an
unbelieving emperor and refused him only in the cause of Christ. Hilary and
Ambrose differ from them, grounding the debt in the possession of the world's
goods. The third asks why true words earn the name of
hypocrisy. Jerome, Chrysostom, Aquinas and Augustine each locate the fault in
the distance between the questioners' words and their hearts, by different
reasoning; Chrysostom, Hilary and Aquinas make Christ's knowledge of the heart
a sign of who he is, while Jerome states the same discernment as a rule for
anyone who answers. Each interpretation gives every appointed element a place
and closes with four senses of its own, literal, allegorical, moral and
anagogical.

\subsection{The appointed elements at a glance}

\begin{maptable}
Entrance Antiphon & \latin{Ego clamavi, quoniam exaudisti me, Deus}; cf.\ Ps
17 (16):6, 8, an adaptation &
One who has called on God because God hears asks him to bend his ear, then to
keep him as the apple of the eye and shelter him under the shadow of his
wings.\\
Collect & \latin{Omnipotens sempiterne Deus, fac nos tibi semper} &
That the assembly may always hold its will devoted to God and serve his
majesty with an undivided heart.\\
First Reading & Is 45:1, 4--6; vv.~2--3 are not read &
The Lord addresses Cyrus as his anointed, grasps his right hand to subdue
nations, calls him by name for Israel's sake though Cyrus does not know him,
and girds him, that all from east to west may know there is no God besides
the Lord.\\
Responsorial Psalm & Ps 96 (95):1 and 3, 4--5, 7--8, 9--10a and 10c; response
from v.~7b &
A new song for all the earth; the Lord's glory declared among the nations;
the gods of the nations are nothing beside him; the families of the peoples
bring him glory and gifts; the Lord reigns and judges the peoples with
justice.\\
Second Reading & 1~Thes 1:1--5b &
Paul, Silvanus and Timothy greet the church of the Thessalonians in God the
Father and the Lord Jesus Christ; they remember its work of faith, labour of
charity and enduring hope, and know its election, because the gospel came in
power and in the Holy Spirit, not in word only.\\
Alleluia verse & Phil 2:15d, 16a &
Those to whom Paul writes shine like lights in the world, holding forth the
word of life.\\
Gospel & Mt 22:15--21 &
The Pharisees' disciples and the Herodians flatter Jesus and ask whether
tribute may lawfully be paid to Caesar; he calls them hypocrites, has them
show the coin and name its image, and commands that Caesar be given what is
Caesar's and God what is God's.\\
Prayer over the Offerings & \latin{Tribue nos, Domine, quaesumus, donis tuis} &
That those who serve God's gifts may do so with a free mind, and that his
purifying grace may cleanse them through the very mysteries in which they
minister.\\
Communion Antiphon, first option & \latin{Ecce oculi Domini super timentes
eum}; cf.\ Ps 33 (32):18--19 &
The eyes of the Lord rest on those who fear him and hope in his mercy, to
rescue them from death and feed them in famine.\\
Communion Antiphon, second option & \latin{Filius hominis venit}; Mk 10:45,
adapted &
The Son of Man came to give his life as a ransom for many.\\
Prayer after Communion & \latin{Fac nos, quaesumus, Domine, caelestium rerum} &
That frequent sharing in heavenly things may make the assembly advance, aided
by temporal benefits and instructed by eternal ones.\\
\end{maptable}

\noindent The two Communion antiphons are alternatives: the Missal prints the
second under \latin{Vel}, and one of them is used. Every other element is
required. The formulary has no Sequence, Offertory antiphon, proper Preface,
proper insert in the Eucharistic Prayer, prayer over the people or solemn
blessing, and the Lectionary gives no other form of any reading. A claim below
that rests on one Communion antiphon holds only where that antiphon is chosen;
a claim that rests on the wording of an antiphon does not hold where another
approved chant is sung in its place.
